% !TeX program = lualatex-dev
% ===================================================================
% mwe_list_pagebreak_bug.tex — reprodução MÍNIMA (sem pacote CRP)
% do bug de \hsize em lista × quebra de página sob phase-III.
%
% Sintoma (camada Leg, 2026-07-11, versão description do \LegArtigo
% a partir do 13º artigo da DUDH): quebra de página no MEIO do
% parágrafo de um item corrompe o \hsize das continuações — cascata
% de \hbox overfull de ~uma palavra por linha, PDF incha ~10×
% (o manual chegou a 3614 páginas).
%
% DIAGNÓSTICO REFINADO (bissecção 2026-07-13): o gatilho NÃO é
% \list + quebra em si — é a invocação do ambiente POR CSNAME
% (\description…\enddescription, como dentro de uma macro), que pula
% os env hooks de que o tagueamento de listas depende. Resultados
% nesta máquina (kernel develop 2026-07-10, tagpdf 1.0b):
%   - este arquivo (csname):          75 págs, 3228 overfull  ← BUG
%   - \begin{description} idêntico:    ~8 págs, 0 overfull    ← ok
%   - sem \DocumentMetadata (csname):  ~8 págs, 0 overfull    ← ok
% \list cru DENTRO de \newenvironment próprio (LegIncisos etc., via
% \begin{…}) é saudável — validado à parte.
%
% Compilar: lualatex-dev mwe_list_pagebreak_bug.tex
% Destino: issue em latex3/tagging-project — rascunho pronto em
% ../../briefings/issue-draft-list-csname-tagging.md; abrir SÓ com
% aprovação de Angelo.
% ===================================================================

\DocumentMetadata{
    lang        = en,
    testphase   = {phase-III},
    tagging     = on,
}

\documentclass{book}

\newcommand\filler{Lorem ipsum dolor sit amet, consectetur adipiscing
elit, sed do eiusmod tempor incididunt ut labore et dolore magna
aliqua. Ut enim ad minim veniam, quis nostrud exercitation ullamco
laboris nisi ut aliquip ex ea commodo consequat. }
\newcommand\longpar{\filler\filler\filler\filler\filler
\filler\filler\filler\filler\filler}

\begin{document}

\chapter*{List item crossing a page break under tagging}

% Gatilho isolado na bissecção de 2026-07-13: a description invocada
% POR CSNAME (\description…\enddescription, forma legítima de macro)
% em vez de \begin{description}…\end{description}. Os env hooks
% (env/description/begin etc.) de que o tagueamento de listas do
% phase-III depende NÃO disparam na forma csname — o estado do
% tagging fica desbalanceado e, quando uma quebra de página cai no
% meio do parágrafo de um item, o \hsize das continuações corrompe
% (cascata de overfull de ~uma palavra por linha; PDF incha ~10×).
% Com \begin{description} idêntico, o documento sai saudável.
\newcount\artn
\artn=1
\loop
    \description
    \item[Article \the\artn.] \filler\filler\filler
    \enddescription
    \advance\artn by 1
\ifnum\artn<31 \repeat

\end{document}
